\documentclass[tikz,border=6pt]{standalone}
\usepackage{ctex}
\usepackage{amsmath}
\usepackage{pgfplots}
\pgfplotsset{compat=1.18}

\begin{document}
\begin{tikzpicture}

\begin{axis}[
  width=11cm, height=8cm,
  xmin=-2.2, xmax=3.5,
  ymin=-0.5, ymax=6.5,
  xlabel={$x$},
  ylabel={$F_t(x) = f(x) - \frac{1}{t}\log(-g(x))$},
  xlabel style={font=\small},
  ylabel style={font=\small, yshift=6pt},
  axis lines=left,
  tick style={font=\small},
  xtick={-2,-1,0,1,2,3},
  ytick={0,1,2,3,4,5,6},
  clip=true,
  legend pos=north west,
  legend style={font=\small, cells={anchor=west}},
]

% 约束边界在 x = 2（即 g(x) = x-2 <= 0, 可行域 x < 2）
% f(x) = 0.5*(x-3)^2（目标函数，最优解在 x=3，但约束要求 x<2，故约束最优在 x=2）
% 障碍函数：phi(x) = -log(2-x)（x < 2）
% F_t(x) = f(x) + (1/t)*phi(x) = 0.5*(x-3)^2 - (1/t)*log(2-x)

% t=0.5 时的 F_t(x)（强障碍，远离边界）
\addplot[
  blue, thick,
  domain=-2.0:1.95, samples=200,
] {0.5*(x-3)^2 + 2*(-ln(2-x))};
\addlegendentry{$t=0.5$（强障碍）}

% t=2 时的 F_t(x)
\addplot[
  orange!80!black, thick,
  domain=-2.0:1.97, samples=200,
] {0.5*(x-3)^2 + 0.5*(-ln(2-x))};
\addlegendentry{$t=2$}

% t=10 时的 F_t(x)（弱障碍，接近原问题）
\addplot[
  green!60!black, thick,
  domain=-2.0:1.99, samples=200,
] {0.5*(x-3)^2 + 0.1*(-ln(2-x))};
\addlegendentry{$t=10$（弱障碍）}

% 原目标函数 f(x) = 0.5*(x-3)^2（无约束）
\addplot[
  gray, dashed, thick,
  domain=-2.0:3.5, samples=100,
] {0.5*(x-3)^2};
\addlegendentry{$f(x)$（无障碍）}

% 约束边界 x=2
\draw[red!80!black, very thick, dashed]
  (axis cs:2, -0.5) -- (axis cs:2, 6.5);
\node[font=\small, red!80!black, anchor=south, align=center]
  at (axis cs:2.0, 5.6) {约束\\边界};

% "屏障墙"标注
\draw[{Stealth[length=5pt]}-{Stealth[length=5pt]}, red!60]
  (axis cs:1.7, 4.8) -- (axis cs:2.0, 4.8);
\node[font=\tiny, red!60, anchor=east] at (axis cs:1.65, 4.8)
  {障碍\\屏障};

% 无约束最优解（x=3，不可行）
\fill[gray!70] (axis cs:3, 0) circle (3pt);
\node[font=\footnotesize, gray!70, anchor=north]
  at (axis cs:3, -0.2) {$x^*_{\text{无约束}}=3$};

% 约束最优解（x*=2，约束边界）
\fill[black] (axis cs:2, 0.5) circle (3pt);
\node[font=\footnotesize, anchor=north west]
  at (axis cs:2.02, 0.4) {$x^*=2$};

% t→∞时各F_t的最小值点（垂直虚线标注）
% t=0.5: 最小值约在 x=0.5 左右
\fill[blue!80] (axis cs:0.2, 3.4) circle (2pt);
\node[font=\tiny, blue!80, anchor=south]
  at (axis cs:0.2, 3.5) {$x^*(0.5)$};

% t=2: 最小值约在 x=1.3
\fill[orange!80!black] (axis cs:1.3, 1.6) circle (2pt);
\node[font=\tiny, orange!80!black, anchor=south]
  at (axis cs:1.3, 1.7) {$x^*(2)$};

% t=10: 最小值约在 x=1.8
\fill[green!60!black] (axis cs:1.78, 0.65) circle (2pt);
\node[font=\tiny, green!60!black, anchor=south]
  at (axis cs:1.78, 0.75) {$x^*(10)$};

% 说明文字
\node[font=\footnotesize, align=center, draw=gray!50, rounded corners, fill=white!80]
  at (axis cs:-0.8, 5.5) {
    $t$ 越大，障碍权重越小\\
    $x^*(t)$ 越靠近约束边界 $x^*=2$
  };

\end{axis}
\end{tikzpicture}
\end{document}
